\textbf{Indústria e crescimento econômico do Rio de Janeiro (2002-2021):
características setoriais e intensidade da desindustrialização
fluminense}

Resumo: O trabalho analisa o processo de desindustrialização do estado
do Rio de Janeiro no período de 2002 até 2021. Se configura como um
trabalho de análise histórico comparativa baseado nas taxas de
crescimento real dos setores industriais e do PIB do Rio de Janeiro e
Brasil. Testa a hipótese de que o estado do Rio de Janeiro passou por um
processo de desindustrialização mais acelerado que outros estados
brasileiros. Como fonte de dados utiliza os dados das Contas Regionais
dos estados brasileiros, disponibilizadas pelo IBGE em parceria com as
secretarias estaduais. As taxas de crescimento real do valor adicionado
da atividade econômica são tratadas como \emph{proxies} de processos de
industrialização e desindustrialização. Como conclusão, o trabalho
apresenta uma proposta de periodização para o processo de
desindustrialização do Rio de Janeiro e corrobora a hipótese de que este
está ocorrendo em ritmo acelerado em relação ao resto do Brasil.

Palavras-chave: indústria, rio de janeiro, desindustrialização, método
comparativo, sociologia histórica

Abstract: The article looks into the deindustrialization process of the
Rio de Janeiro (RJ) state, in Brazil, between 2002 and 2021. It's a
historical-comparative sociology based on the real growth rates of RJ's
and Brazil's industry segments. It tests the hypothesis that RJ is
undergoing a faster deindustrialization process than the other Brazilian
states. The main source of data are the regional accounts, made
available by Brazilian Institute of Geography and Statistics (IBGE) and
local government agencies. Real growth rates of the added value to GDP
are used as proxies of industrialization and deindustrialization. The
author's conclusion corroborates the hypothesis. It also proposes a
periodization of RJ's deindustrialization.

Keywords: industry, rio de janeiro, deindustrialization, comparative
method, historical sociology

\textbf{Introdução}

O trabalho aqui apresentado é uma sociologia histórico-comparativa do
processo de desindustrialização recente (2002-2021) do estado do Rio de
Janeiro. É baseado na análise das séries temporais das taxas de
crescimento real dos setores industriais. As séries são comparadas entre
si e com os dados gerais do Brasil. O objetivo é oferecer uma descrição
mais acurada do ritmo e periodização do processo de desindustrialização
recente do estado, oferecendo insumos tanto para análises causais, como
para análises sócio-históricas do processo de reestruturação econômica
em andamento. A hipótese é de que o estado do Rio de Janeiro passou por
um processo de desindustrialização recente mais intenso do que o Brasil
e outros estados.

O debate sobre a desindustrialização brasileira tem ganhado proeminência
nos últimos anos. Evidência disso é a formulação de políticas públicas
específicas para lidar com a questão
\href{https://www.zotero.org/google-docs/?wo4Op5}{(CNDI e MDIC, 2024)}.
Isso pois, para além de um fenômeno a ser compreendido, o processo de
desindustrialização tem sido problematizado enquanto prejudicial ao
potencial de crescimento econômico
\href{https://www.zotero.org/google-docs/?jIz1rX}{(Oreiro e Feijó,
2010)}. Esse cenário é distinto da avaliação de parte da teoria
econômica que em grande maneira considerava esse processo como esperado
ou mesmo positivo
\href{https://www.zotero.org/google-docs/?z9I4ma}{(Oreiro e Marconi,
2012)}. Essa preocupação com a atualidade do processo coloca o debate
brasileiro sobre desindustrialização em trajetória distinta de boa parte
da literatura internacional, que enfatiza um processo que considera em
larga medida encerrado
\href{https://www.zotero.org/google-docs/?ta9ftR}{(High, 2013; Lawson,
2020; Strangleman, Rhodes e Linkon, 2013)}, mesmo que com consequências
perenes \href{https://www.zotero.org/google-docs/?MUfXCU}{(Linkon, 2018;
Mah, 2012)}.

Dentro dos campos das ciências sociais há, entretanto, uma problemática
adicional. Se a desindustrialização tem sido concebida como um processo
histórico de impacto social significativo, que encerra uma era
industrial
\href{https://www.zotero.org/google-docs/?ky4Ed4}{(Strangleman, 2017)},
é preciso entender quais as características desse processo transformam
seus possíveis resultados. Nesse ponto, uma das características mais
importantes dos processos de mudança social é seu ritmo e duração
\href{https://www.zotero.org/google-docs/?G2kQTr}{(Abbott, 1984;
Pierson, 2004)}. Mudanças sociais bruscas ou cumulativas, em longa
duração, têm chamado a atenção de cientistas sociais e historiadores que
buscam dar a devida atenção à dimensão temporal dos processos
históricos.

\textbf{História, ritmo e mudança social}

Skocpol \href{https://www.zotero.org/google-docs/?N5g7Uw}{(1984)}
identifica que a sociologia histórica tinha predileção por pensar
``grandes questões'' e em lógica comparativa. Da mesma maneira, Tilly
\href{https://www.zotero.org/google-docs/?H0ADAb}{(1984)} enfatiza as
grandes comparações como particularidade de certa sociologia preocupada
com processos de transformação em escala macro. Fortemente influenciada
pelo campo da história, essa sociologia discutirá a importância de
processos de mudança em longa duração e perspectivas sistêmicas
globais\href{https://www.zotero.org/google-docs/?4OGNLt}{(Smith, 1991)}.

Adicionalmente, como Kreuzer
\href{https://www.zotero.org/google-docs/?Ic6cdI}{(2023)}argumenta,
processos históricos precisam ser pensados para além da dimensão do
tempo histórico -- o ordenamento e periodização de eventos e datas --
abarcando também concepções de tempo físico. Isso significa inserir uma
dimensão de ritmo onde, além da sequencialidade e ordem, a velocidade e
duração dos processos é relevante para compreensão de seus resultados. O
uso de séries temporais é particularmente útil para esses fins, não
sendo aqui visto como um substituto, mas sim um complemento a análise de
narrativa densa típica dos estudos historiográficos. Em específico, o
uso de variáveis que capturem uma dimensão de ritmo dos processos
históricos é útil para identificação de dois padrões: primeiramente,
ajuda a identificar processos de mudança lenta e gradual,
particularmente aquelas que envolvem algum quesito de dependência de
trajetória; em segundo lugar, permite a identificação de inflexões nas
tendências, ajudando a percepção do que a literatura passou a denominar
``conjunturas críticas''
\href{https://www.zotero.org/google-docs/?YTbI1n}{(Mahoney, 2000, 2004;
Pierson, 2015; Thelen, 2010)}.

Dentre essas ``grandes'' questões que mais movimentaram historiadores e
cientistas sociais está a emblemática discussão sobre
``desenvolvimento''. Desenvolvimento é uma metáfora historiográfica, mas
que carrega significativo impacto político e cognitivo
\href{https://www.zotero.org/google-docs/?4FCLf9}{(Fernández Sebastián,
2024)}. Como Santos argumenta, para além de categoria nativa, é também
``capaz de conferir legitimidade e justificação a agentes políticos e
econômicos'' \href{https://www.zotero.org/google-docs/?tkmOPq}{(Santos,
2016, p. 295)}.

No caso brasileiro ela foi de particular importância nas discussões
sobre a natureza ``dependente''
\href{https://www.zotero.org/google-docs/?GvmQDz}{(Evans, 2018)} desse
possível desenvolvimento nacional, discussão que em grande parte se
debruçava sobre o papel da industrialização. Mas a participação da
indústria na economia brasileira foi elemento central não apenas para
historiadores e cientistas sociais brasileiros, mas para a classe
política e setores das elites
\href{https://www.zotero.org/google-docs/?JiVg8Y}{(Barbosa, 2020)}. Dado
esse papel, a desindustrialização acaba por ser mais do que uma mudança
econômica, é uma transformação simultânea dos arranjos institucionais e
políticos nacionais e do que se concebe como visão de país. Não à toa,
surgimento de projetos de \emph{neoindustrialização,} como o Nova
Indústria Brasil \href{https://www.zotero.org/google-docs/?ZxpRYL}{(CNDI
e MDIC, 2024)}, ou de ``reconversão industrial'', como no caso do
projeto de produção dos caças Gripen no ABC paulista
\href{https://www.zotero.org/google-docs/?t3Oo3o}{(Ramalho e Conceição,
2024)}, acabam condicionados a transformações nas conjunturas políticas
nacionais e locais.

\textbf{Desindustrialização do Brasil e do estado do Rio de Janeiro}

A categoria indústria é por vezes mobilizada de forma distinta na
literatura sobre processos de desindustrialização, com ênfase desigual
sobre seus setores. Enquanto o processo histórico de desindustrialização
é associado em grande parte a atividades extrativas -- como nos casos
inglês \href{https://www.zotero.org/google-docs/?dvLJw6}{(Tomlinson,
2016)} e alemão
\href{https://www.zotero.org/google-docs/?TzOvAP}{(Berger, Wicke e
Golombek, 2017)} --, seus estudos mais recentes enfatizam a indústria de
transformação, em grande parte ancorada na perspectiva de uma posição
privilegiada desta no processo de crescimento econômico. Aqui propõe-se
uma análise dos quatro setores nativamente considerados como parte da
indústria \href{https://www.zotero.org/google-docs/?iSA9Xh}{(CNI,
{[}s.d.{]})}: 1-indústria de transformação, 2-extrativa, 3-construção e
4-serviços de água, esgoto, gás e energia (doravante AEGE) -- com
significativa ênfase nos três primeiros.

Já a desindustrialização é tipicamente lida como uma redução relativa do
valor adicionado da indústria no PIB e uma redução de sua participação
no estoque de empregos
\href{https://www.zotero.org/google-docs/?k039Yz}{(Tregenna, 2009)}.
Suas causas estão relacionadas a processos de mudança tecnológica,
possibilitando o crescimento da produtividade industrial acima do
crescimento da demanda
\href{https://www.zotero.org/google-docs/?niRiWF}{(Rowthorn e Ramaswamy,
1999)} e a realocação produtiva de unidades fabris para territórios de
custo mais baixo
\href{https://www.zotero.org/google-docs/?Qt10Wg}{(Alderson, 1999)}.
Nesse sentido, a desindustrialização é um fenômeno cuja identificação
exige um recorte territorial para sua identificação ou não, podendo
ocorrer simultaneamente a processos de industrialização em outros
territórios \href{https://www.zotero.org/google-docs/?dQKhXi}{(Bluestone
e Harrison, 1982)}. Entre seus impactos, além dos custos mais imediatos
das perdas de empregos para comunidades e famílias afetadas, está um
contínuo processo de declínio dos territórios afetados, que
\href{https://www.zotero.org/google-docs/?tIvUMy}{(Mah, 2012)} denomina
``ruinação''. Linkon
\href{https://www.zotero.org/google-docs/?MfNH7j}{(2018)} vai observar
também como os processos de desindustrialização tem impacto nas
populações trabalhadoras mesmo após o período mais intenso de
fechamentos industriais..

Um conjunto de autores pontua uma relação entre PIB per capita e
desindustrialização. Nessa tese, a partir de certo ponto o crescimento
do PIB per capita a demanda por serviços cresceria em ritmo mais
acelerado que a demanda por bens industriais
\href{https://www.zotero.org/google-docs/?uvhkcx}{(Rodrik, 2016)}. Em
contraste, Rodrik
\href{https://www.zotero.org/google-docs/?KF8pD8}{(2016)} observa como
processos de desindustrialização passaram a ocorrer em países com renda
per capita abaixo do antigo patamar estimado como causador do fenômeno.
A esse fenômeno ele denominou ``desindustrialização prematura'',
tipicamente ocorrendo em países com processos de industrialização
tardia\footnote{Para discussão sobre industrialização e industrialização
  tardia, ver Mello (1982)}.

É esse processo ``prematuro'' que Oreiro e Feijó
\href{https://www.zotero.org/google-docs/?hrJ4AN}{(2010)} identificam
ocorrer no Brasil, em contraponto ao processo de desindustrialização
naturalizado como decorrência esperada do processo de desenvolvimento,
que muitas vezes a teoria econômica apresenta. Em que pese a
sobrestimação da participação da indústria de transformação no PIB
brasileiro antes dos anos 2000 em função de mudanças metodológicas do
sistema de contas nacionais, como mostra Morceiro
\href{https://www.zotero.org/google-docs/?yaPvp0}{(2021)}, há forte
evidência da ocorrência de um processo de desindustrialização no Brasil
entre meados dos anos 1980 e fins dos ano 1990
\href{https://www.zotero.org/google-docs/?EpdASN}{(Oreiro e Marconi,
2012)}. Dúvidas surgiram em princípios dos anos 2000 se esse fenômeno
seria continuado, uma mudança estrutural da economia brasileira, ou
apenas um fenômeno pontual dos anos 80, como defendeu Nassif
\href{https://www.zotero.org/google-docs/?Lcq024}{(2008; 2015)}.
Entretanto, evidências mais recentes apontam solidez do processo de
desindustrialização ao longo das duas primeiras décadas do século XXI,
como pode ser visto em Morceiro e Guilhoto
\href{https://www.zotero.org/google-docs/?GHtpZP}{(2020)}, Maia
\href{https://www.zotero.org/google-docs/?aTAxBm}{(2020)} e Feijó,
Lamônica e Lima
\href{https://www.zotero.org/google-docs/?1SorMi}{(2019)}.

No caso do Rio de Janeiro o processo de desindustrialização está também
ligado a características próprias do estado. Sobral
\href{https://www.zotero.org/google-docs/?zV6TKb}{(2016)} aponta
relações de um ``vazio'' produtivo na região metropolitana com uma
concentração da atividade econômica em torno da metrópole. Em outro
artigo, o mesmo autor aponta as consequências econômicas negativas desse
processo de reestruturação econômica do estado, já perceptíveis no
início dos anos 2000
\href{https://www.zotero.org/google-docs/?cnRPuc}{(Sobral, 2017)}.
Compartilhando a preocupação com a concentração e desigualdade econômica
vista da perspectiva espacial, Vianna da Cruz
\href{https://www.zotero.org/google-docs/?gibX7I}{(2013)} discute os
impactos da problemática industrial para o desenvolvimento do estado e
da excessiva dependência deste no setor petrolífero. No quesito emprego,
Souza e Veríssimo
\href{https://www.zotero.org/google-docs/?oj5qly}{(2019)} apontam que,
ao menos até meados dos anos de 2010, se melhor que o cenário nacional,
o Rio de Janeiro já enfrentava decréscimo de cerca de 3,23\% do emprego
industrial na participação do emprego total -- saindo de 13,2\% em 1996
para 10,22\% em 2014, portanto, quase ¼ da participação no estoque de
emprego.

Entre as causas do caso nacional, há particular ênfase na literatura
sobre o papel das taxas de câmbio, particularmente a ideia de que esta
teria estado sobrevalorizada no período
\href{https://www.zotero.org/google-docs/?SmfxoM}{(Marconi e Rocha,
2012; Oreiro e Feijó, 2010)}. Diretamente ligada à questão cambial há
também frequente identificação de uma ``doença holandesa''
\href{https://www.zotero.org/google-docs/?3ti55W}{(Cunha, Lelis e
Fligenspan, 2013)}ou ``a desindustrialização causada pela apreciação da
taxa real de câmbio resultante da descoberta de recursos naturais
escassos'' \href{https://www.zotero.org/google-docs/?rWMrUB}{(Oreiro e
Feijó, 2010, p. 222)}. Isso implica numa relação direta entre a
desindustrialização do setor de transformação e a performance do setor
extrativo.

Referente aos indicadores utilizados, a maioria dos estudos enfatiza a
participação do valor adicionado da indústria -- na maioria dos casos de
transformação -- no PIB. Porém, para além dos já citados problemas
metodológicos referentes a como o dado é contabilizado, também há
questionamentos sobre a adequação deste indicador devido ao excessivo
impacto que a inflação setorial tem na participação de cada setor
\href{https://www.zotero.org/google-docs/?NYfMT1}{(Torres e Cavalieri,
2015)}. Em outras palavras, um dado setor industrial pode aumentar
significativamente sua participação no PIB sem nenhum aumento produtivo,
basta a inflação setorial ser superior a do PIB. Como resposta a essas
questões, propõe-se na próxima seção do artigo o uso de taxas de
crescimento dos setores econômicos em lugar de sua participação no PIB.

\textbf{Métodos}

O recorte temporal aqui analisado vai de 2002 até 2021. Dado que o
cálculo da taxa de crescimento exige a comparação entre dois anos têm-se
dados do crescimento real de 2003 até 2021. O recorte foi escolhido em
função da limitação dos dados das contas regionais, sendo assim restrito
a série mais atual. Utilizar dados anteriores foi descartado por
questões de tempo e espaço, já que os dados exigiram tratamento
particular em função de mudanças metodológicas pontuadas acima e um
espaço de discussão adicional não comportado no limite deste artigo.

São utilizados os dados provenientes das contas regionais,
disponibilizados pelo IBGE e elaborados em parceria com as secretarias
estaduais \href{https://www.zotero.org/google-docs/?uwDtJh}{(IBGE,
{[}s.d.{]})}. Foi utilizado o dado setorial de cada estado para o
cálculo da taxa de crescimento anual dos setores da indústria e do PIB
de 2003 até 2021, utilizando-se sempre o Valor Adicionado (VA) do ano
corrente aos preços do ano anterior e o VA do ano anterior a preços
correntes.

\(\frac{VALOR\ ADICIONADO\ do\ ANO\ CORRENTE\ em\ PREÇOS\ DO\ ANO\ ANTERIOR}{VALOR\ ANO\ ANTERIOR\ em\ PREÇOS\ CORRENTES}\  = \ Taxa\ de\ Crescimento\ (1)\ \)

Adicionalmente, apenas para o caso do Rio de Janeiro foi calculada a
participação dos setores industriais no PIB para o período de 2002 até
2021 de forma a demonstrar o peso dos setores industriais na economia
fluminense e assim contextualizar o impacto distinto das taxas de
crescimento destes na economia. Os cálculos utilizaram o VA de cada
setor da indústria dividido pelo VA do total das atividades (i.e. PIB).
O cálculo da categoria Indústria foi feito somando o VA dos quatro
setores industriais e dividindo-se pelo VA total das atividades.

\(\frac{VALOR\ ADICIONADO\ do\ SETOR\ no\ ANO\ CORRENTE}{VALOR\ ADICIONADO\ das\ ATIVIDADES\ \ TOTAIS\ no\ ANO\ CORRENTE}\  = \ Participação\ do\ setor\ no\ PIB\ do\ RJ\ (2)\)

\(Transformação\  + Extrativa\  + \ Construção\  + AEGE\  = \ Indústria\ Geral\ (3)\)

O uso das taxas de crescimento como \emph{proxy} da desindustrialização
captura melhor o sua aceleração e desaceleração, além de permitir uma
comparação entre setores industriais que podem ter participação distinta
na atividade econômica entre Rio, Brasil e outros estados. Taxas de
crescimento mais baixas dos setores industriais em relação ao PIB
estadual indicariam um processo de desindustrialização e a diferença,
sua intensidade.

Por fim, os dados foram analisados a partir de técnicas de análise de
séries temporais, com ênfase nas ciências sociais
\href{https://www.zotero.org/google-docs/?fQMQx8}{(Shin, 2017)}. Na
prática, realizou-se a análise gráfica das séries, a comparação de
similaridade entre ela e a comparação das médias de crescimento para
recortes do período. No primeiro caso, suavizou-se a série temporal via
média móvel de três anos como forma de facilitar a visualização da
tendência da série\footnote{A média móvel foi utilizada \textbf{apenas
  para fins de visualização}, não para cálculo das médias de crescimento
  do período ou correlação entre as séries temporais. Os erros relativos
  à média móvel para cada série encontram-se no anexo I ao fim do
  artigo.}. Foi estabelecido um número mínimo de dois anos para cálculo
de cada média, o que na prática eliminou os dados de 2003, assim, os
gráficos apresentados têm seu eixo X de 2004 até 2021. No segundo caso
utilizou-se a correlação de Pearson para mensurar a similaridade entre
séries temporais distintas, cujo resultado que equivale a uma correlação
cruzada com \emph{lag=0}
\href{https://www.zotero.org/google-docs/?au1UC7}{(Shin, 2017)}.
Entende-se que uma correlação alta implica que ambas as séries temporais
possam compartilhar fatores causais, ao passo que baixas correlações
aumentam a probabilidade de fatores causais distintos terem prioridade
em cada uma das séries. Os recortes do período foram analisados a partir
da identificação de suas similaridades e diferenças com objetivo de
estabelecer ou eliminar possíveis causas suficientes e necessárias, em
linha com abordagens histórico-comparativas
\href{https://www.zotero.org/google-docs/?9E83JL}{(Mahoney, 2004)}.

\textbf{Resultados}

Primeiramente contextualiza-se o peso relativo dos setores industriais
no PIB fluminense. O gráfico 1 (abaixo) demonstra como a indústria
extrativa tem grande participação neste, com fortes flutuações. Assim,
ao considerar-se as taxas de crescimento é importante ter em mente o
peso deste setor. Em contrapartida, o setor de transformação parece ter
um declínio mais constante. Já o setor de construção parecia ter
assumido papel mais significativo a partir de meados da primeira década
dos anos 2000, apenas para retroceder em meados da seguinte. O setor
relativo à AEGE representa parte bem menor do PIB. Em função disso e
também dos limites de espaço e tempo, ele não será enfocado no restante
da análise.

Gráfico 1 - Participação dos setores industriais no PIB do estado do RJ

\includegraphics[width=5.16667in,height=3.74581in]{media/image2.png}

Fonte: autoria própria a partir das Contas Regionais (IBGE)

Já um olhar sobre as taxas de crescimento real para o período (gráfico
2, p.9) esclarece em parte as razões das flutuações. No caso da
indústria extrativa existe uma gradual redução das taxas de crescimento,
inclusive com resultados abaixo do PIB na primeira década do século, com
recuperação na segunda, o que a priori não coincide com seu crescimento
(e queda) de participação no PIB (gráfico 1, abaixo). O período de
performance mais fraco coincide entretanto com um período de forte
valorização do preço das \emph{commodities,} associado a boa performance
da economia durante a primeira década do século
\href{https://www.zotero.org/google-docs/?KSrQl7}{(Fonseca, Arend e
Guerrero, 2020)}, e consequente crescimento de sua participação no PIB
estadual. Assim, o crescimento da participação do setor extrativo no PIB
parece dever-se não a boa performance e investimentos, mas a mudanças
nos preços. Coincide, entretanto, com um período relativamente longo de
valorização cambial, no qual o dólar passou de 3,0783 reais em 2003,
para 1,6750 reais em 2011
\href{https://www.zotero.org/google-docs/?76BKq1}{(IPEA, 2024)}.

Gráfico 2 - Média Móvel de 3 anos das Taxas de crescimento real do
estado do RJ

\includegraphics[width=5.26563in,height=3.70388in]{media/image8.png}

Fonte: autoria própria a partir das Contas Regionais (IBGE)

Já o setor de transformação exibe constante fragilidade, com taxas de
crescimento real abaixo do PIB fluminense em quase todo o período
(gráfico 2, acima). Isso não implica necessariamente decrescimento,
apesar desse ter sido frequente, mas sim uma desindustrialização
relativa à economia estadual. É importante notar também que a
fragilidade da indústria de transformação frente ao PIB fluminense é
mais significativa quando considerado que o PIB do estado como um todo
cresceu abaixo do nacional e estava entre os piores resultados dos
estados (gráfico 3, p.10).

Gráfico 3 - Média Móvel de 3 anos das taxas de crescimento dos PIB
estaduais e nacional

\includegraphics[width=5.86458in,height=4.2415in]{media/image6.png}

Fonte: autoria própria a partir das Contas Regionais (IBGE)

Considerando-se as taxas de crescimento do PIB carioca e também os
movimentos dos setores industriais ao nível nacional, propõe-se a
periodização desse processo de desindustrialização em três etapas. De
2002 até 2008 houve um período de ``desindustrialização relativa'', ou
seja, taxas de crescimento dos setores industriais foram positivas,
porém no geral inferiores ao PIB estadual. Um segundo período de 2009
até 2014 é marcado por piora das taxas industriais (com parcial exceção
do setor de construção) e sua entrada em territórios negativos. Já de
2015 até 2021 tem-se primeiro a agudização da crise, especialmente entre
2015 e 2016, com gradual recuperação das taxas de crescimento até o fim
do período. As tabelas 1 e 2 abaixo sintetizam os resultados dos setores
industriais para os três períodos.

Tabela 1 - Média de crescimento real do PIB e setores industriais do RJ
nos períodos

\begin{longtable}[]{@{}
  >{\raggedright\arraybackslash}p{(\linewidth - 8\tabcolsep) * \real{0.3967}}
  >{\raggedright\arraybackslash}p{(\linewidth - 8\tabcolsep) * \real{0.1476}}
  >{\raggedright\arraybackslash}p{(\linewidth - 8\tabcolsep) * \real{0.1476}}
  >{\raggedright\arraybackslash}p{(\linewidth - 8\tabcolsep) * \real{0.1568}}
  >{\raggedright\arraybackslash}p{(\linewidth - 8\tabcolsep) * \real{0.1513}}@{}}
\toprule\noalign{}
\begin{minipage}[b]{\linewidth}\raggedright
\textbf{Setor/Período}
\end{minipage} & \begin{minipage}[b]{\linewidth}\raggedright
\textbf{2003-2008}
\end{minipage} & \begin{minipage}[b]{\linewidth}\raggedright
\textbf{2009-2014}
\end{minipage} & \begin{minipage}[b]{\linewidth}\raggedright
\textbf{2015-2021}
\end{minipage} & \begin{minipage}[b]{\linewidth}\raggedright
\textbf{2003-2021}
\end{minipage} \\
\begin{minipage}[b]{\linewidth}\raggedright
PIB RJ
\end{minipage} & \begin{minipage}[b]{\linewidth}\raggedright
2.48\%
\end{minipage} & \begin{minipage}[b]{\linewidth}\raggedright
2.13\%
\end{minipage} & \begin{minipage}[b]{\linewidth}\raggedright
-0.44\%
\end{minipage} & \begin{minipage}[b]{\linewidth}\raggedright
1.19\%
\end{minipage} \\
\begin{minipage}[b]{\linewidth}\raggedright
Transformaçao Rio
\end{minipage} & \begin{minipage}[b]{\linewidth}\raggedright
0.73\%
\end{minipage} & \begin{minipage}[b]{\linewidth}\raggedright
-0.39\%
\end{minipage} & \begin{minipage}[b]{\linewidth}\raggedright
-1.32\%
\end{minipage} & \begin{minipage}[b]{\linewidth}\raggedright
-0.82\%
\end{minipage} \\
\begin{minipage}[b]{\linewidth}\raggedright
Extrativa Rio
\end{minipage} & \begin{minipage}[b]{\linewidth}\raggedright
2.43\%
\end{minipage} & \begin{minipage}[b]{\linewidth}\raggedright
-0.72\%
\end{minipage} & \begin{minipage}[b]{\linewidth}\raggedright
6.08\%
\end{minipage} & \begin{minipage}[b]{\linewidth}\raggedright
2.71\%
\end{minipage} \\
\begin{minipage}[b]{\linewidth}\raggedright
Construção RIO
\end{minipage} & \begin{minipage}[b]{\linewidth}\raggedright
1.35\%
\end{minipage} & \begin{minipage}[b]{\linewidth}\raggedright
5.39\%
\end{minipage} & \begin{minipage}[b]{\linewidth}\raggedright
-5.81\%
\end{minipage} & \begin{minipage}[b]{\linewidth}\raggedright
-0.12\%
\end{minipage} \\
\midrule\noalign{}
\endhead
\bottomrule\noalign{}
\endlastfoot
\end{longtable}

Fonte: autoria própria a partir das Contas Regionais (IBGE)

Tabela 2 - Diferença de crescimento dos setores industriais (RJ) para o
nível nacional e PIB (RJ)

\begin{longtable}[]{@{}l@{}}
\toprule\noalign{}
\endhead
\bottomrule\noalign{}
\endlastfoot
\end{longtable}

Fonte: autoria própria a partir das Contas Regionais (IBGE)

Cada um dos períodos propostos acima, entretanto, para além das taxas de
crescimento, está imerso em variáveis importantes do ponto de vista
histórico, especialmente mudanças políticas e econômicas na escala
nacional e internacional. Assim, o período de 2002 até 2008 é marcado
pela primeira eleição de Luiz Inácio Lula da Silva (2003-2010) à
presidência, que, a despeito de continuidades na política
macroeconômica, marca também mudanças do ponto de vista da política de
valorização do salário mínimo
\href{https://www.zotero.org/google-docs/?mo5q5w}{(Cosenza, 2022)}. O
período se encerra em 2008, com estabelecimento de forte crise econômica
internacional.

O novo período, de 2009 até 2014 se inicia com duas inflexões:
primeiramente o ano de 2009 será marcado por significativa desaceleração
econômica no Brasil (gráfico 3 acima). Em grande parte em função desta
haverá uma mudança significativa da política econômica. De particular
importância é o lançamento do primeiro Programa de Aceleração do
Crescimento (PAC) em 2007
\href{https://www.zotero.org/google-docs/?2jrHE4}{(Rousseff, 2007)} e
subsequentemente o lançamento do Minha Casa Minha Vida, em 2009, este em
larga medida um programa de oferta de crédito subsidiado para aquisição
de moradias de baixa renda
\href{https://www.zotero.org/google-docs/?lpgneb}{(BRASIL, 2009)}. Essas
mudanças atravessaram o final do segundo governo Lula e a presidência de
Dilma Rousseff (2011-2016).

O final do período é marcado por significativas transformações
econômicas e políticas. Particularmente, o ano de 2013 vê emergência de
manifestações sociais de rua, conjugado com o início das atividades da
Operação Lava Jato. Já em 2014, e parcialmente relacionado a esses
fatores \href{https://www.zotero.org/google-docs/?baledF}{(DIEESE,
2021)}, tem se o início de uma desaceleração econômica significativa, em
grande parte vinculada a uma forte queda no preço das
\emph{commodities}. Esta foi iniciado em meados de 2011 e intensificada
entre 2014 e 2015
\href{https://www.zotero.org/google-docs/?sPw1rR}{(Fonseca, Arend e
Guerrero, 2020)}.

Por fim, o período de 2015 até 2021 é marcado pela intensificação da
crise econômica -- em parte já em curso em 2014. Cobre tanto o período
de impeachment da ex-presidente Dilma Rousseff, como os governos Michel
Temer (2016-2018), Jair Bolsonaro (2019-2022) e a primeira parte da
crise do coronavírus (2020-2022 aprox.). Também em 2015 têm-se
significativas mudanças na política econômica a nível federal, com
adoção de uma política fiscal de redução de gastos
\href{https://www.zotero.org/google-docs/?uUV1Vw}{(Fonseca, Arend e
Guerrero, 2020)}.

Já quanto aos recortes setoriais para cada um dos períodos e começando
pelo setor de transformação, é notável que este tem piores taxas de
crescimento em relação ao dado nacional (tabela 2, p.11). Olhando
especificamente para os períodos de 2002 até 2008 e 2015 até 2021 o Rio
de Janeiro registrou algumas das piores taxas para o setor entre todos
os estados. Ao mesmo tempo, se de 2009 até 2014 a taxa de crescimento
foi inferior ao dos outros dois períodos, importante notar que foi esse
o momento onde a performance da indústria de transformação do Rio de
Janeiro mais se aproximou do dado nacional (tabela 2, p.10). Isso parece
indicar que, para além do cenário macroeconômico, existem outros
limitantes estruturais ao crescimento da indústria de transformação
fluminense, que mesmo em momentos de ``recuperação'' exibe baixo
crescimento.

Considerando a variável câmbio, é relevante como, após contínua
valorização, o real passa a se desvalorizar de 1,765 em 2011 até 5,395
em 2021 \href{https://www.zotero.org/google-docs/?PaAtIJ}{(IPEA, 2024)}.
Essa mudança não parece, entretanto, sustentar efeitos longos para a
indústria de transformação, a qual, após breve aceleração até 2012,
perde força e passa a ter taxas de crescimento inferiores ao período
anterior (gráfico 4 acima).

Analisando a correlação (tabela 3, p.13) do setor de transformação do
Rio de Janeiro com o do Brasil e o PIB fluminense, os resultados indicam
maior correlação com o primeiro. Isso sugere que o setor de
transformação da indústria fluminense foi mais afetado por dinâmicas do
setor de transformação ao nível nacional do que com a performance
econômica do estado. Dito isso, ainda há correlação razoável (0.64, com
R2 de 0.41) entre a taxa de crescimento real da indústria de
transformação do Rio de Janeiro e o PIB do estado.

Gráfico 4 - Média Móvel de 3 anos da taxa de crescimento real da
indústria de transformação

\includegraphics[width=6.35417in,height=4.25814in]{media/image5.png}

Fonte: autoria própria a partir das Contas Regionais (IBGE)

Tabela 3 - Correlação de PEARSON entre as séries temporais de taxas de
crescimento real

\begin{longtable}[]{@{}
  >{\raggedright\arraybackslash}p{(\linewidth - 8\tabcolsep) * \real{0.2320}}
  >{\raggedright\arraybackslash}p{(\linewidth - 8\tabcolsep) * \real{0.1680}}
  >{\raggedright\arraybackslash}p{(\linewidth - 8\tabcolsep) * \real{0.2100}}
  >{\raggedright\arraybackslash}p{(\linewidth - 8\tabcolsep) * \real{0.1900}}
  >{\raggedright\arraybackslash}p{(\linewidth - 8\tabcolsep) * \real{0.2000}}@{}}
\toprule\noalign{}
\begin{minipage}[b]{\linewidth}\raggedright
Correlação*
\end{minipage} &
\multicolumn{2}{>{\raggedright\arraybackslash}p{(\linewidth - 8\tabcolsep) * \real{0.3780} + 2\tabcolsep}}{%
\begin{minipage}[b]{\linewidth}\raggedright
\textbf{R}
\end{minipage}} &
\multicolumn{2}{>{\raggedright\arraybackslash}p{(\linewidth - 8\tabcolsep) * \real{0.3900} + 2\tabcolsep}@{}}{%
\begin{minipage}[b]{\linewidth}\raggedright
\textbf{R2}
\end{minipage}} \\
\begin{minipage}[b]{\linewidth}\raggedright
RIO (set\textbackslash or)
\end{minipage} & \begin{minipage}[b]{\linewidth}\raggedright
Brasil (setor)
\end{minipage} & \begin{minipage}[b]{\linewidth}\raggedright
PIB (RJ)
\end{minipage} & \begin{minipage}[b]{\linewidth}\raggedright
Brasil (setor)
\end{minipage} & \begin{minipage}[b]{\linewidth}\raggedright
PIB (RJ)
\end{minipage} \\
\begin{minipage}[b]{\linewidth}\raggedright
\textbf{Transformação}
\end{minipage} & \begin{minipage}[b]{\linewidth}\raggedright
0.7812
\end{minipage} & \begin{minipage}[b]{\linewidth}\raggedright
0.6409
\end{minipage} & \begin{minipage}[b]{\linewidth}\raggedright
0.6102
\end{minipage} & \begin{minipage}[b]{\linewidth}\raggedright
0.4107
\end{minipage} \\
\begin{minipage}[b]{\linewidth}\raggedright
\textbf{Extrativo}
\end{minipage} & \begin{minipage}[b]{\linewidth}\raggedright
-0.1025
\end{minipage} & \begin{minipage}[b]{\linewidth}\raggedright
-0.1960
\end{minipage} & \begin{minipage}[b]{\linewidth}\raggedright
0.0105
\end{minipage} & \begin{minipage}[b]{\linewidth}\raggedright
0.0384
\end{minipage} \\
\begin{minipage}[b]{\linewidth}\raggedright
\textbf{Construção}
\end{minipage} & \begin{minipage}[b]{\linewidth}\raggedright
0.9313
\end{minipage} & \begin{minipage}[b]{\linewidth}\raggedright
0.7974
\end{minipage} & \begin{minipage}[b]{\linewidth}\raggedright
0.8673
\end{minipage} & \begin{minipage}[b]{\linewidth}\raggedright
0.6358
\end{minipage} \\
\multirow{3}{=}{\begin{minipage}[b]{\linewidth}\raggedright
\emph{P-Values}

\emph{(respectivamente)}
\end{minipage}} & \begin{minipage}[b]{\linewidth}\raggedright
7.87E-05
\end{minipage} & \begin{minipage}[b]{\linewidth}\raggedright
3.11E-03
\end{minipage} &
\multicolumn{2}{>{\raggedright\arraybackslash}p{(\linewidth - 8\tabcolsep) * \real{0.3900} + 2\tabcolsep}@{}}{%
\multirow{3}{=}{\begin{minipage}[b]{\linewidth}\raggedright
\end{minipage}}} \\
& \begin{minipage}[b]{\linewidth}\raggedright
6.76E-01
\end{minipage} & \begin{minipage}[b]{\linewidth}\raggedright
4.21E-01
\end{minipage} \\
& \begin{minipage}[b]{\linewidth}\raggedright
7.14E-09
\end{minipage} & \begin{minipage}[b]{\linewidth}\raggedright
4.34E-05
\end{minipage} \\
\midrule\noalign{}
\endhead
\bottomrule\noalign{}
\endlastfoot
\end{longtable}

Fonte: autoria própria a partir das Contas Regionais (IBGE)

A indústria extrativa do Rio de Janeiro também exibiu performance
inferior à média nacional em quase todo o período (gráfico 5 abaixo).
Essa tendência só se modifica no período de 2015 a 2021, sendo capaz de
reverter a tendência de desindustrialização setorial e estabelecer um
crescimento médio 1,52 ponto percentual absoluto acima do PIB do estado
para todo o período (tabela 2, p.11). No mesmo gráfico 5, vê-se como a
taxa de crescimento da indústria extrativa inicia próxima da taxa do PIB
fluminense, mas vai gradualmente perdendo força, em parte em
concomitância com a queda cambial.

A participação desta no PIB (gráfico 1, p.8) continua alta em virtude
dos preços elevados das commodities, mas a taxa real de crescimento
mostra outro cenário. É justamente durante o período de inversão da
tendência de valorização do câmbio no qual a indústria extrativa voltará
a ver suas taxas de crescimento acelerarem (gráfico 6, abaixo). Isso é
em especial curioso quando se observa que essa aceleração se deu
justamente no período da chamada crise da \emph{commodities} (entre 2014
e 2015). Esses dados parecem contrastar, ao menos para o caso do Rio de
Janeiro, com a hipótese de que o investimento na indústria extrativa,
motivado por um câmbio valorizado e altos preços de \emph{commodities},
seria responsável pela desindustrialização do setor de transformação.
Isso não implica que o câmbio não exerça influência, entretanto, o
período aproximadamente de 2011 até 2015, de forte valorização cambial,
foi onde a indústria de transformação fluminense teve desempenho mais
próximo ao pib estadual (gráfico 4, p.13), em contraposição ao período
de 2003 até 2010, de câmbio em forte valorização. Dito isso, não é causa
suficiente -- e nem mesmo necessária -- para uma performance negativa do
setor em relação ao PIB do estado, como o período de 2015 até 2020
mostra.

Na análise da correlação da série temporal com o mesmo setor ao nível
nacional e com o PIB do estado (tabela 3, p.13), em ambos os casos houve
correlação negativa. Tal resultado pode indicar uma dinâmica bastante
independente da indústria extrativa do estado do Rio de Janeiro.
Entretanto, é importante frisar que foram os únicos casos onde P
\textgreater{} 0.05 (0.68 e 0.42 respectivamente), o que reduz a
confiança no dado. Com fins de verificar se poderia, para a indústria
extrativa, existir um \emph{lag} entre as séries, foi feita a correlação
cruzada, mas como pode ser visto no correlograma (gráfico 7 abaixo), não
parece haver qualquer lag onde a correlação seja significativamente
maior.

Gráfico 5 - Média Móvel de 3 anos da taxa de crescimento real da
indústria extrativa

\includegraphics[width=5.09318in,height=3.65523in]{media/image7.png}

Fonte: autoria própria a partir das Contas Regionais (IBGE)

Gráfico 6 - Média móvel (3 anos) da Ind. extrativa e de transformação
versus câmbio

\includegraphics[width=5.33854in,height=3.85434in]{media/image3.png}

Fonte: autoria própria a partir das Contas Regionais (IBGE) e Ipea

Grafico 7 - Correlograma da indústria extrativa com PIB e setor nacional

\includegraphics[width=4.79688in,height=3.51164in]{media/image1.png}

Fonte: autoria própria a partir das Contas Regionais (IBGE)

Por fim, o setor de construção (gráfico 8 abaixo) mostra o mais
acentuado ciclo de aceleração e desaceleração das taxas de crescimento
(como fica nítido no gráfico 2, p.9). De particular importância parece
ser a implementação do primeiro PAC, com particular peso para o programa
Minha Casa Minha Vida, a partir de 2009. Esse também é o período onde o
setor de construção teve crescimento acima do PIB, revertendo em parte a
desindustrialização setorial.

Gráfico 8 - Média Móvel de 3 anos da taxa de crescimento real da
indústria de construção

\includegraphics[width=6.01042in,height=4.26436in]{media/image4.png}

Fonte: autoria própria a partir das Contas Regionais (IBGE)

Em contraposição, o setor experimenta forte contração desde 2013. Uma
hipótese explicativa é o impacto da Operação Lava Jato nas atividades
das principais construtoras e obras públicas. Esse cenário parece se
intensificar após a contração econômica de 2015, exibindo gradual
recuperação a partir de 2017. O período a partir de 2017 também mostra a
aceleração da taxa de crescimento do setor de construção fluminense para
acima do nível nacional, em contraposição com a maior parte da série
temporal. De todo modo, a indústria de construção ainda continuou tendo
taxas de crescimento muito inferiores ao PIB estadual, mantendo o
processo de desindustrialização (tabela 2, p.11).

Analisando a correlação percebe-se que o desempenho do setor de
construção é o mais fortemente correlacionado tanto ao PIB fluminense
como ao setor no nível nacional (tabela 3, p.13). Particularmente, a
correlação de 0.93 com o PIB ao nível nacional pode indicar uma
prevalência de fatores nacionais na performance do setor de construção
fluminense. Essa hipótese dialoga com a relação proposta acima entre o
setor de construção e a intensificação --ou retração -- de políticas
públicas voltadas à realização de obras ou subsídios para o setor
habitacional.

\textbf{Discussão e conclusões}

O artigo corroborou a hipótese de que o Rio de Janeiro tem passado por
um processo de desindustrialização mais acelerado que os demais estados
brasileiros. Essa desindustrialização, pode ser vista de duas
perspectivas. Pode ser considerada como uma desindustrialização em
relação à atividade econômica do estado. Nesse caso, tanto indústria de
transformação como a de construção viram média de crescimento inferior
ao PIB estadual no período (perda de 2.01 e 1.32 ponto percentual
absoluto, respectivamente; tabela 2, p.11), significando uma
desindustrialização ``clássica''. A indústria extrativa cresceu em
relação ao PIB (1.52 ponto percentual absoluto), mas esse crescimento
esteve quase todo concentrado no período mais recente (2015-2021). Já em
relação ao Brasil podemos falar de uma segunda forma
``desindustrialização'', representada pela perda de importância do Rio
de Janeiro para a atividade industrial nacional. Sinal disso, indústrias
de transformação e construção tiveram taxas de crescimento real
inferiores à taxa em nível nacional (-1.49 e 1.66 ponto percentual
absoluto, respectivamente; tabela 2, p.11). Já nesse recorte a indústria
extrativa teve estabilidade (ganho de 0.09 ponto percentual absoluto;
tabela 2, p.11).

Há, portanto, claras distinções setoriais. Os testes de correlação entre
as séries temporais parecem indicar que os setores extrativo, de
transformação e construção estão vinculados em grau distinto a
condicionantes no nível nacional (câmbio, crescimento do PIB, políticas
públicas federais, conjuntura política, etc). A performance do setor
extrativo, particularmente, parece muito mais independente da
performance do mesmo setor em nível nacional do que o setor de
construção, estando a indústria de transformação numa posição
intermediária. Isso indica um maior peso de fatores regionais ou
especificamente setoriais em determinadas indústrias do que em outras.
De toda maneira, fica claro que existem elementos no caso do Rio de
Janeiro que contribuem para uma performance abaixo de outros estados
brasileiros na maioria dos casos e períodos.

Do ponto de vista das causas explicativas, é necessário complexificar,
ao menos para o caso fluminense, a hipótese de reprimarização
relacionada à sobrevalorização cambial como explicação central para o
fenômeno da desindustrialização. Isso pois não parece haver
concomitância da desvalorização cambial com uma melhora contínua na
indústria de transformação (gráfico 6, p.16). Da mesma forma, a
indústria extrativa não performou de forma inferior durante o choque de
preços das \emph{commoditie}s, ao contrário, viu suas taxas de
crescimento acelerarem (gráfico 5, p.14). Evidentemente não se descarta
o impacto do câmbio no processo de desindustrialização -- ou a
possibilidade de reprimarização --, mas esta não parece ser causa
suficiente e nem necessária para a desindustrialização do Rio de
Janeiro.

Do ponto de vista do impacto das políticas públicas -- e da retirada
destas -- a indústria de construção parece o mais claro exemplo deste.
Os dados indicam grande efeito do primeiro PAC e do Minha Casa Minha
Vida sobre as taxas de crescimento do setor. Ao mesmo tempo, a
paralisação de grande parte das atividades em correlação com a Operação
Lava Jato mostra o forte peso do fato político nos ciclos de aceleração
e desaceleração da desindustrialização fluminense. Vale citar que,
apesar da Operação Lava Jato também ter enfocado empresas do setor
extrativo, em particular a Petrobrás, com fortes efeitos nas estratégias
de investimento desta
\href{https://www.zotero.org/google-docs/?X39guU}{(DIEESE, 2021)}, a
mesma desaceleração do crescimento real não se refletiu na indústria
extrativa. Isso não elimina a possibilidade de impacto desta no setor
extrativo, mas indica que esse impacto não foi o mais determinante para
as taxas de crescimento do setor nesse período. Do ponto de vista
teórico, a implicação parece ser de que o efeito da conjuntura política
pode afetar em diferentes graus setores industriais distintos.

Por fim, vale pontuar limitações deste estudo e direções para sua
continuidade. Primeiramente, o recorte temporal pode ser estendido, ao
menos até o início dos anos 1980. Além disso, não foram feitas
regressões multivariadas, o que não permite modelar ou detectar o
impacto de diversas variáveis nas taxas de crescimento, tais como
inflação e investimento. Em terceiro lugar, não foram utilizados dados
de produção física e emprego, que certamente auxiliarão a ter um retrato
mais completo do fenômeno e suas particularidades setoriais.

Quanto à análise histórica dos períodos -- e a periodização proposta --,
esta é claramente preliminar e enfoca a escala nacional. Analisar com
mais profundidade como as sequências de eventos políticos e econômicos
influenciaram os setores industriais e mapear o conjunto de políticas
públicas estaduais mais relevantes é importante continuidade dessa
pesquisa. Ainda importante é considerar que a escolha dos casos
comparativos pode ser sofisticada, realizando comparação com dados de
estados específicos, considerados casos mais próximos ou mais distantes
ao do Rio de Janeiro (demograficamente, em sua composição setorial,
processo histórico, etc). Dito isso, se o retrato aqui apresentado do
processo de desindustrialização fluminense auxiliar continuidades dessa
investigação, indicando períodos relevantes para observação aprofundada
e a importância de levar em conta as diferenças setoriais, o objetivo
terá sido atingido.

Anexo 1 - erros da média móvel

\begin{longtable}[]{@{}
  >{\raggedright\arraybackslash}p{(\linewidth - 6\tabcolsep) * \real{0.3518}}
  >{\raggedright\arraybackslash}p{(\linewidth - 6\tabcolsep) * \real{0.2348}}
  >{\raggedright\arraybackslash}p{(\linewidth - 6\tabcolsep) * \real{0.2258}}
  >{\raggedright\arraybackslash}p{(\linewidth - 6\tabcolsep) * \real{0.1875}}@{}}
\toprule\noalign{}
\begin{minipage}[b]{\linewidth}\raggedright
\textbf{Série/Erro da MM}
\end{minipage} & \begin{minipage}[b]{\linewidth}\raggedright
\textbf{mae}
\end{minipage} & \begin{minipage}[b]{\linewidth}\raggedright
\textbf{mape}
\end{minipage} & \begin{minipage}[b]{\linewidth}\raggedright
\textbf{rmse}
\end{minipage} \\
\begin{minipage}[b]{\linewidth}\raggedright
PIB Rio de Janeiro
\end{minipage} & \begin{minipage}[b]{\linewidth}\raggedright
0,01252
\end{minipage} & \begin{minipage}[b]{\linewidth}\raggedright
0,62024
\end{minipage} & \begin{minipage}[b]{\linewidth}\raggedright
0,00025
\end{minipage} \\
\begin{minipage}[b]{\linewidth}\raggedright
PIB Brasil
\end{minipage} & \begin{minipage}[b]{\linewidth}\raggedright
0,01738
\end{minipage} & \begin{minipage}[b]{\linewidth}\raggedright
2,72567
\end{minipage} & \begin{minipage}[b]{\linewidth}\raggedright
0,00047
\end{minipage} \\
\begin{minipage}[b]{\linewidth}\raggedright
Transf. Rio de Janeiro
\end{minipage} & \begin{minipage}[b]{\linewidth}\raggedright
0,03878
\end{minipage} & \begin{minipage}[b]{\linewidth}\raggedright
8,41341
\end{minipage} & \begin{minipage}[b]{\linewidth}\raggedright
0,00249
\end{minipage} \\
\begin{minipage}[b]{\linewidth}\raggedright
Transf. Brasil
\end{minipage} & \begin{minipage}[b]{\linewidth}\raggedright
0,03598
\end{minipage} & \begin{minipage}[b]{\linewidth}\raggedright
1,15073
\end{minipage} & \begin{minipage}[b]{\linewidth}\raggedright
0,00184
\end{minipage} \\
\begin{minipage}[b]{\linewidth}\raggedright
Extrat. Rio de Janeiro
\end{minipage} & \begin{minipage}[b]{\linewidth}\raggedright
0,03982
\end{minipage} & \begin{minipage}[b]{\linewidth}\raggedright
1,15272
\end{minipage} & \begin{minipage}[b]{\linewidth}\raggedright
0,00247
\end{minipage} \\
\begin{minipage}[b]{\linewidth}\raggedright
Extrat. Brasil
\end{minipage} & \begin{minipage}[b]{\linewidth}\raggedright
0,03914
\end{minipage} & \begin{minipage}[b]{\linewidth}\raggedright
1,65896
\end{minipage} & \begin{minipage}[b]{\linewidth}\raggedright
0,00219
\end{minipage} \\
\begin{minipage}[b]{\linewidth}\raggedright
Const. Rio de Janeiro
\end{minipage} & \begin{minipage}[b]{\linewidth}\raggedright
0,04848
\end{minipage} & \begin{minipage}[b]{\linewidth}\raggedright
1,29798
\end{minipage} & \begin{minipage}[b]{\linewidth}\raggedright
0,00311
\end{minipage} \\
\begin{minipage}[b]{\linewidth}\raggedright
Const. Brasil
\end{minipage} & \begin{minipage}[b]{\linewidth}\raggedright
0,03687
\end{minipage} & \begin{minipage}[b]{\linewidth}\raggedright
1,30759
\end{minipage} & \begin{minipage}[b]{\linewidth}\raggedright
0,00219
\end{minipage} \\
\midrule\noalign{}
\endhead
\bottomrule\noalign{}
\endlastfoot
\end{longtable}

Fonte: autoria própria a partir de IPEA

Bibliografía

\href{https://www.zotero.org/google-docs/?p3HTgZ}{ABBOTT, A. Event
Sequence and Event Duration: Colligation and Measurement.
\textbf{Historical Methods: A Journal of Quantitative and
Interdisciplinary History}, v. 17, n. 4, p. 192--204, out. 1984.}

\href{https://www.zotero.org/google-docs/?p3HTgZ}{ALDERSON, A. S.
Explaining Deindustrialization: Globalization, Failure, or Success?
\textbf{American Sociological Review}, v. 64, n. 5, p. 701--721, out.
1999.}

\href{https://www.zotero.org/google-docs/?p3HTgZ}{BARBOSA, A. D. F.
``Developmentalist Brazil'' (1945-1964) as a concept: historicizing and
(re)periodizing development in Brazil. \textbf{Brazilian Journal of
Political Economy}, v. 40, n. 2, p. 332--354, jun. 2020.}

\href{https://www.zotero.org/google-docs/?p3HTgZ}{BERGER, S.; WICKE, C.;
GOLOMBEK, J. Burdens of Eternity? \textbf{The Public Historian}, v. 39,
n. 4, p. 21--43, 1 nov. 2017.}

\href{https://www.zotero.org/google-docs/?p3HTgZ}{BLUESTONE, B.;
HARRISON, B. \textbf{The deindustrialization of America: plant closings,
community abandonment, and the dismantling of basic industry}. New York:
Basic Books, 1982.}

\href{https://www.zotero.org/google-docs/?p3HTgZ}{BRASIL. \textbf{LEI
N\textsuperscript{o}} \textbf{11.977, DE 7 DE JULHO DE 2009}, 2009.
Disponível em:
\textless https://www.planalto.gov.br/ccivil\_03/\_ato2007-2010/2009/lei/l11977.htm\textgreater.
Acesso em: 10 set. 2024}

\href{https://www.zotero.org/google-docs/?p3HTgZ}{CNDI, C. N. DE D. I.;
MDIC, M. DO D., Indústria, Comércio e Serviços. \textbf{Nova indústria
Brasil -- Nova indústria Brasil -- forte, transformadora e
sustentável\,: Plano de Ação para a neoindustrialização 2024-2026}.
Brasília: CNDI, MDIC, 2024. Disponível em:
\textless https://www.gov.br/mdic/pt-br/composicao/se/cndi/plano-de-acao/nova-industria-brasil-plano-de-acao.pdf\textgreater.
Acesso em: 29 mar. 2024.}

\href{https://www.zotero.org/google-docs/?p3HTgZ}{CNI. \textbf{CNI --
Perfil da Indústria Brasileira}. Disponível em:
\textless https://industriabrasileira.portaldaindustria.com.br/\textgreater.
Acesso em: 2 maio. 2024.}

\href{https://www.zotero.org/google-docs/?p3HTgZ}{COSENZA, A. C. Um
brasil dos trabalhadores? Aspectos do desempenho econômico brasileiro
entre 2003 e 2016. \emph{Em}: SOUZA, L. E. S. DE; PREVIDELLI, M. DE F.
S. DO C. (Eds.). . \textbf{História econômica do Brasil contemporâneo}.
Coleção Novos estudos de história econômica do Brasil. Niterói, RJ\,:
São Paulo, SP: Eduff\,; HUCITEC Editora, 2022. .}

\href{https://www.zotero.org/google-docs/?p3HTgZ}{CUNHA, A. M.; LELIS,
M. T. C.; FLIGENSPAN, F. B. Desindustrialização e comércio exterior:
evidências recentes para o Brasil. \textbf{Brazilian Journal of
Political Economy}, v. 33, n. 3, p. 463--485, set. 2013.}

\href{https://www.zotero.org/google-docs/?p3HTgZ}{DIEESE.
\textbf{IMPLICAÇÕES ECONÔMICAS INTERSETORIAIS DA OPERAÇÃO LAVA JATO},
2021.}

\href{https://www.zotero.org/google-docs/?p3HTgZ}{EVANS, P. B.
\textbf{Dependent Development: the Alliance of Multinational, State, and
Local Capital in Brazil}. Princeton, NJ: Princeton University Press,
2018.}

\href{https://www.zotero.org/google-docs/?p3HTgZ}{FEIJÓ, C.; LAMÔNICA,
M. T.; LIMA, S. S. Financial liberalization and structural change: the
Brazilian case in the 2000s. \textbf{Economia e Sociedade}, v. 28, n. 1,
p. 177--200, 2019.}

\href{https://www.zotero.org/google-docs/?p3HTgZ}{FERNÁNDEZ SEBASTIÁN,
J. \textbf{Key metaphors for history}. London\,; New York, NY:
Routledge, Taylor \& Francis Group, 2024.}

\href{https://www.zotero.org/google-docs/?p3HTgZ}{FONSECA, P. C. D.;
AREND, M.; GUERRERO, G. A. Política econômica, instituições e classes
sociais: os governos do Partido dos Trabalhadores no Brasil.
\textbf{Economia e sociedade (São Paulo, Brazil)}, v. 29, n. 3, p.
779--809, 2020.}

\href{https://www.zotero.org/google-docs/?p3HTgZ}{HIGH, S. ``The Wounds
of Class'': A Historiographical Reflection on the Study of
Deindustrialization, 1973--2013. \textbf{History Compass}, v. 11, n. 11,
p. 994--1007, nov. 2013.}

\href{https://www.zotero.org/google-docs/?p3HTgZ}{IBGE. \textbf{Contas
Regionais do Brasil \textbar{} IBGE}. Disponível em:
\textless https://www.ibge.gov.br/estatisticas/economicas/contas-nacionais/9054-contas-regionais-do-brasil.html\textgreater.
Acesso em: 2 maio. 2024.}

\href{https://www.zotero.org/google-docs/?p3HTgZ}{IPEA. \textbf{Taxa de
câmbio}. Disponível em:
\textless http://www.ipeadata.gov.br/ExibeSerie.aspx?serid=31924\textgreater.
Acesso em: 6 set. 2024.}

\href{https://www.zotero.org/google-docs/?p3HTgZ}{KREUZER, M.
\textbf{The Grammar of Time: A Toolbox for Comparative Historical
Analysis}. 1. ed. {[}s.l.{]} Cambridge University Press, 2023.}

\href{https://www.zotero.org/google-docs/?p3HTgZ}{LAWSON, C. Making
sense of the ruins: The historiography of deindustrialisation and its
continued relevance in neoliberal times. \textbf{History Compass}, v.
18, n. 8, p. e12619, ago. 2020.}

\href{https://www.zotero.org/google-docs/?p3HTgZ}{LINKON, S. L.
\textbf{The half-life of deindustrialization: working-class writing
about economic restructuring}. Ann Arbor: University of Michigan Press,
2018.}

\href{https://www.zotero.org/google-docs/?p3HTgZ}{MAH, A.
\textbf{Industrial ruination, community, and place: landscapes and
legacies of urban decline}. Toronto: University of Toronto Press, 2012.}

\href{https://www.zotero.org/google-docs/?p3HTgZ}{MAHONEY, J. Path
Dependence in Historical Sociology. \textbf{Theory and Society}, v. 29,
n. 4, p. 507--548, 2000.}

\href{https://www.zotero.org/google-docs/?p3HTgZ}{\_\_\_.
Comparative-Historical Methodology. \textbf{Annual Review of Sociology},
v. 3, p. 81--101, 2004.}

\href{https://www.zotero.org/google-docs/?p3HTgZ}{MAIA, B. A. DE A. Há
desindustrialização no Brasil? Um estudo da abordagem clássica e de
análises alternativas entre 1998 e 2014. \textbf{Economia e sociedade
(São Paulo, Brazil)}, v. 29, n. 2, p. 549--579, 2020.}

\href{https://www.zotero.org/google-docs/?p3HTgZ}{MARCONI, N.; ROCHA, M.
Taxa de câmbio, comércio exterior e desindustrialização precoce: o caso
brasileiro. \textbf{Economia e Sociedade}, v. 21, n. spe, p. 853--888,
dez. 2012.}

\href{https://www.zotero.org/google-docs/?p3HTgZ}{MORCEIRO, P. C.
Influência metodológica na desindustrialização brasileira.
\textbf{Brazilian Journal of Political Economy}, v. 41, n. 4, p.
700--722, 2021.}

\href{https://www.zotero.org/google-docs/?p3HTgZ}{MORCEIRO, P. C.;
GUILHOTO, J. J. M. Adensamento produtivo e esgarçamento do tecido
industrial brasileiro. \textbf{Economia e Sociedade}, v. 29, n. 3, p.
835--860, 2020.}

\href{https://www.zotero.org/google-docs/?p3HTgZ}{NASSIF, A. Há
evidências de desindustrialização no Brasil? \textbf{Brazilian Journal
of Political Economy}, v. 28, n. 1, p. 72--96, mar. 2008.}

\href{https://www.zotero.org/google-docs/?p3HTgZ}{NASSIF, L.; TEIXEIRA,
L.; ROCHA, F. Houve redução do impacto da indústria na economia
brasileira no período 1996-2009? Uma análise das matrizes
insumo-produto. \textbf{Economia e Sociedade}, v. 24, n. 2, p. 355--378,
2015.}

\href{https://www.zotero.org/google-docs/?p3HTgZ}{OREIRO, J. L.; FEIJÓ,
C. A. Desindustrialização: conceituação, causas, efeitos e o caso
brasileiro. \textbf{Revista de Economia Política}, v. 30, n. 2, p.
219--232, jun. 2010.}

\href{https://www.zotero.org/google-docs/?p3HTgZ}{OREIRO, J. L.;
MARCONI, N. Teses Equivocadas no Debate sobre Desindustrialização e
Perda de Competitividade da Indústria Brasileira. 2012.}

\href{https://www.zotero.org/google-docs/?p3HTgZ}{PIERSON, P.
\textbf{Politics in time: history, institutions, and social analysis}.
Princeton: Princeton University Press, 2004.}

\href{https://www.zotero.org/google-docs/?p3HTgZ}{\_\_\_. Power and Path
Dependence. \emph{Em}: MAHONEY, J.; THELEN, K. (Eds.). .
\textbf{Advances in Comparative-Historical Analysis}. Cambridge:
Cambridge Univ. Press, 2015. .}

\href{https://www.zotero.org/google-docs/?p3HTgZ}{RAMALHO, J. R.;
CONCEIÇÃO, J. J. DA. \textbf{Caças supersônicos e o ABC Paulista:
Tecnologia e reconversão industrial}. São Paulo, SP: Editora Papagaio,
2024.}

\href{https://www.zotero.org/google-docs/?p3HTgZ}{RODRIK, D. Premature
deindustrialization. \textbf{Journal of Economic Growth}, v. 21, n. 1,
p. 1--33, mar. 2016.}

\href{https://www.zotero.org/google-docs/?p3HTgZ}{ROUSSEFF, D. PROGRAMA
DE ACELERAÇÃO DO CRESCIMENTO. 2007.}

\href{https://www.zotero.org/google-docs/?p3HTgZ}{ROWTHORN, R.;
RAMASWAMY, R. Growth, Trade, and Deindustrialization. \textbf{IMF Staff
Papers}, v. 46, n. 1, p. 18--41, 1999.}

\href{https://www.zotero.org/google-docs/?p3HTgZ}{SANTOS, R. S. P. DOS.
Desenvolvimento econômico e mudança social: a Vale e a mineração na
Amazônia Oriental. \textbf{Caderno CRH}, v. 29, n. 77, p. 295--312, ago.
2016.}

\href{https://www.zotero.org/google-docs/?p3HTgZ}{SHIN, Y. \textbf{Time
series analysis in the social sciences: the fundamentals}. Oakland,
California: University of California Press, 2017.}

\href{https://www.zotero.org/google-docs/?p3HTgZ}{SKOCPOL, T. (ED.).
\textbf{Vision and method in historical sociology}. Cambridge
{[}Cambridgeshire{]}\,; New York: Cambridge University Press, 1984.}

\href{https://www.zotero.org/google-docs/?p3HTgZ}{SMITH, D. \textbf{The
rise of historical sociology}. Philadelphia: Temple University Press,
1991.}

\href{https://www.zotero.org/google-docs/?p3HTgZ}{SOBRAL, B. L. B.
Desindustrialização e questão metropolitana: o caso da ``arrebentação
urbana'' na periferia da região metropolitana do Rio de Janeiro.
\textbf{Geosul}, v. 31, n. 62, p. 193--220, 2016.}

\href{https://www.zotero.org/google-docs/?p3HTgZ}{\_\_\_. A falácia da
``inflexão econômica positiva'': algumas características da
desindustrialização fluminense e do ``vazio produtivo'' em sua periferia
metropolitana. \textbf{Cadernos do Desenvolvimento Fluminense}, n. 10,
2017.}

\href{https://www.zotero.org/google-docs/?p3HTgZ}{SOUZA, I. E. L. DE;
VERÍSSIMO, M. P. Produção e emprego industrial nos estados brasileiros:
evidências de desindustrialização. \textbf{Nova economia (Belo
Horizonte, Brazil)}, v. 29, n. 1, p. 75, 2019.}

\href{https://www.zotero.org/google-docs/?p3HTgZ}{STRANGLEMAN, T.
Deindustrialisation and the Historical Sociological Imagination: Making
Sense of Work and Industrial Change. \textbf{Sociology}, v. 51, n. 2, p.
466--482, abr. 2017.}

\href{https://www.zotero.org/google-docs/?p3HTgZ}{STRANGLEMAN, T.;
RHODES, J.; LINKON, S. Introduction to Crumbling Cultures:
Deindustrialization, Class, and Memory. \textbf{International Labor and
Working-Class History}, v. 84, p. 7--22, 2013.}

\href{https://www.zotero.org/google-docs/?p3HTgZ}{THELEN, K. Beyond
Comparative Statics: Historical Institutional Approaches to Stability
and Change In the Political Economy of Labor. \emph{Em}: MORGAN, G.
\emph{et al.} (Eds.). . \textbf{The Oxford Handbook of Comparative
Institutional Analysis}. 1. ed. {[}s.l.{]} Oxford University Press,
2010. p. 41--62.}

\href{https://www.zotero.org/google-docs/?p3HTgZ}{TILLY, C. \textbf{Big
structures, large processes, huge comparisons}. New York: Russell Sage
Foundation, 1984.}

\href{https://www.zotero.org/google-docs/?p3HTgZ}{TOMLINSON, J.
De-industrialization Not Decline: A New Meta-narrative for Post-war
British History. \textbf{Twentieth Century British History}, v. 27, n.
1, p. 76--99, mar. 2016.}

\href{https://www.zotero.org/google-docs/?p3HTgZ}{TORRES, R. L.;
CAVALIERI, H. Uma crítica aos indicadores usuais de desindustrialização
no Brasil. \textbf{Revista de economia política}, v. 35, n. 4, p.
859--877, 2015.}

\href{https://www.zotero.org/google-docs/?p3HTgZ}{TREGENNA, F.
Characterising deindustrialisation: An analysis of changes in
manufacturing employment and output internationally. \textbf{Cambridge
Journal of Economics}, v. 33, n. 3, p. 433--466, 1 maio 2009.}

\href{https://www.zotero.org/google-docs/?p3HTgZ}{VIANNA DA CRUZ, J. L.
Os desafios da construção do desenvolvimento no Estado do Rio de
Janeiro. \textbf{Cadernos do Desenvolvimento Fluminense}, v. 0, n. 2, p.
47--80, 10 jul. 2013.}
